\documentclass[12pt,a4paper]{article}
\usepackage[utf8]{inputenc}
\usepackage[T1]{fontenc}
\usepackage[indonesian]{babel}
\usepackage{geometry}
\usepackage{longtable}
\usepackage{array}
\usepackage{booktabs}
\usepackage{enumitem}
\usepackage[hidelinks,breaklinks]{hyperref}
\geometry{margin=2.5cm}

\begin{document}

\begin{center}
\textbf{\Large RENCANA PEMBELAJARAN SEMESTER (RPS)}\\
\textbf{Berbasis Outcome-Based Education (OBE)}\\[0.5cm]
\textbf{MATA KULIAH PRAKTIKUM / LABORATORIUM}\\[0.25cm]
\textbf{PROGRAM STUDI TEKNIK INFORMATIKA}\\
\textbf{FAKULTAS TEKNOLOGI INFORMASI}\\
\textbf{UNIVERSITAS BALE BANDUNG}\\
\end{center}

\vspace{0.75cm}

% ============================================================
% 1. IDENTITAS MATA KULIAH (BAGIAN A — KERANGKA OBE PRAKTIKUM)
% ============================================================
\section*{1. Identitas Mata Kuliah}

\begin{tabular}{ll}
Nama Program Studi & : Teknik Informatika (S1) \\
Nama Mata Kuliah & : Praktikum Data Science \\
Kode Mata Kuliah & : TIF-xxx-prak \\
Semester & : [Ganjil/Genap] --- diselaraskan dengan penjadwalan kurikulum \\
SKS / Bobot Kredit & : 1 SKS (100\% praktikum; komponen teori = 0 SKS) \\
Rincian bentuk & : Praktikum = 1 SKS (setara \textbf{P = 1}) \\
Jumlah Pertemuan & : 16 (enam belas) sesi praktikum efektif \\
Dosen Pengampu & : [Nama Dosen, M.Kom.] \\
Koordinator Praktikum & : [Nama Koordinator, M.Kom.] (opsional) \\
Tim Asisten Lab & : [Nama asisten] (opsional) \\
Tanggal Penyusunan / Revisi & : [Tanggal] / [Revisi ke-...] \\
\end{tabular}

\vspace{0.45cm}

\noindent\textbf{Prasyarat dan ko-requisit:} Mahasiswa telah mengambil atau \textbf{sedang mengikuti secara bersamaan} mata kuliah \textit{Data Science} (atau setara) agar konsep teori (algoritma, metrik, metodologi) dapat dihubungkan dengan praktik di laboratorium.

\vspace{0.35cm}

\noindent\textbf{Deskripsi mata kuliah:} Praktikum Data Science memberikan pengalaman terstruktur dalam menerapkan alur kerja \textit{data science} menggunakan \textbf{Python} dan \textbf{Jupyter Notebook} (alternatif: JupyterLab, notebook di VS Code, atau Google Colab). Mahasiswa membangun \textit{notebook} yang merekam langkah dari pemahaman data, prapemrosesan, analisis eksploratif, pemodelan dengan \texttt{scikit-learn}, evaluasi, hingga penyampaian \textit{insight} sederhana. Mata kuliah mendukung profil lulusan yang mampu mengimplementasikan komputasi dan algoritma pada bahasa pemrograman untuk permasalahan berbasis data.

\vspace{0.35cm}

\noindent\textbf{Beban waktu (acuan nasional \& lokal):} Sesuai standar umum penyusunan beban kredit di perguruan tinggi, \textbf{1 SKS praktikum} lazim setara dengan sekitar \textbf{170 menit kegiatan terstruktur per minggu} selama satu semester untuk komponen praktikum; rincian jam tatap muka di lab dapat dipecah sesuai jadwal prodi (mis.\ 170 menit kontinyu atau dua slot) dengan tetap memenuhi \textbf{ketentuan akademik Universitas Bale Bandung}.

\vspace{0.35cm}

\noindent\textbf{Ruang, perangkat lunak, dan fasilitas wajib:} Laboratorium komputer dengan Python 3.x, Jupyter Notebook/JupyterLab, paket: NumPy, Pandas, Matplotlib, Seaborn, scikit-learn; koneksi internet untuk dokumentasi dan dataset terbuka; penyimpanan untuk mengumpulkan berkas \texttt{.ipynb} dan data. \textbf{K3 dan etika:} penggunaan peralatan listrik yang aman, tidak memakan/minum di dekat perangkat, dan kepatuhan terhadap \textbf{lisensi dataset} serta \textbf{privasi data}.

\vspace{0.5cm}

% ============================================================
% 2. CPL PRODI
% ============================================================
\section*{2. CPL Prodi yang Dibebankan pada Mata Kuliah}

\begin{itemize}[leftmargin=*]
  \item \textbf{CPL01} --- Kemampuan menganalisis persoalan \textit{computing} yang kompleks serta menerapkan prinsip-prinsip \textit{computing} dan disiplin ilmu relevan lainnya untuk mengidentifikasi solusi, dengan mempertimbangkan wawasan perkembangan ilmu transdisiplin.

  \item \textbf{CPL03} --- Kemampuan merancang dan menganalisis algoritma untuk menyelesaikan permasalahan organisasi secara optimal, serta memilih dan menerapkannya pada bahasa pemrograman tertentu.
\end{itemize}

\vspace{0.5cm}

% ============================================================
% 3. CPMK PRAKTIKUM (KEMAMPUAN AKHIR MK) & PEMETAAN KE MK TEORI
% ============================================================
\section*{3. Capaian Pembelajaran Mata Kuliah (CPMK) --- Praktikum}

CPMK berikut adalah kemampuan \textbf{psikomotorik/keterampilan komputasi} di akhir mata kuliah praktikum. Disusun agar \textbf{selaras} dengan CPMK mata kuliah teori \textit{Data Science} (lihat \texttt{silabus\_data\_science.tex}): CPMK-P berfokus pada implementasi di Python/Jupyter.

\begin{itemize}[leftmargin=*]
  \item \textbf{CPMK-P1} --- (CPL01; jembatan ke CPMK1 teori) Mahasiswa mampu menyiapkan lingkungan praktik Python/Jupyter, menerapkan K3 laboratorium dasar, serta merumuskan permasalahan sederhana dan tahapan \textit{CRISP-DM} ringkas untuk eksperimen berbasis data.

  \item \textbf{CPMK-P2} --- (CPL01, CPL03; CPMK2 teori) Mahasiswa mampu memuat data dari sumber tabular, melakukan prapemrosesan (\textit{cleaning}, transformasi, rekayasa fitur dasar), serta mengeksekusi Analisis Data Eksploratif (EDA) dan visualisasi menggunakan Pandas, Matplotlib, dan Seaborn dalam notebook.

  \item \textbf{CPMK-P3} --- (CPL03; CPMK3 teori) Mahasiswa mampu membangun, melatih, dan memprediksi dengan model \textit{machine learning} (regresi, klasifikasi, \textit{ensemble} ringkas, klastering) menggunakan API \texttt{scikit-learn} secara konsisten (\texttt{fit}/\texttt{predict}/\texttt{transform}).

  \item \textbf{CPMK-P4} --- (CPL01, CPL03; CPMK4 teori) Mahasiswa mampu menghitung dan menginterpretasikan metrik evaluasi, menerapkan validasi silang dan penyetelan hiperparameter sederhana, serta menyusun narasi dan visual \textit{insight} di notebook untuk mendukung pengambilan keputusan.
\end{itemize}

\vspace{0.5cm}

% ============================================================
% 4. SUB-CPMK PRAKTIKUM (AGREGAT)
% ============================================================
\section*{4. Sub-CPMK / Indikator Pencapaian (Praktikum)}

Indikator berikut merangkum tahapan keterampilan lab (lebih ringkas dari 14 sub teori pada MK Data Science, tetap konsisten secara substansi).

\begin{itemize}[leftmargin=*]
  \item \textbf{Sub-CPMK-P1:} Mahasiswa mampu mengoperasikan Jupyter Notebook, mengelola kernel, dan menjelaskan etika serta lisensi penggunaan data pada tugas praktikum.

  \item \textbf{Sub-CPMK-P2:} Mahasiswa mampu memuat data (CSV/Excel, \textit{upload} lokal) dan menginspeksi struktur, tipe, serta kualitas awal dengan Pandas.

  \item \textbf{Sub-CPMK-P3:} Mahasiswa mampu menangani \textit{missing values}, duplikasi, \textit{outlier} dasar, serta melakukan \textit{encoding}, \textit{scaling}, dan rekayasa fitur sederhana.

  \item \textbf{Sub-CPMK-P4:} Mahasiswa mampu melaksanakan EDA (univariat/bivariat) dan memvisualisasikan pola dengan Matplotlib/Seaborn.

  \item \textbf{Sub-CPMK-P5:} Mahasiswa mampu menjelaskan perbedaan \textit{supervised} dan \textit{unsupervised}, memecah data latih/uji, dan menggunakan estimators \texttt{scikit-learn} secara benar.

  \item \textbf{Sub-CPMK-P6:} Mahasiswa mampu mengimplementasikan regresi (mis.\ regresi linear) dan menafsirkan kesalahan prediksi.

  \item \textbf{Sub-CPMK-P7:} Mahasiswa mampu mengimplementasikan klasifikasi dasar (mis.\ regresi logistik, K-NN, pohon keputusan).

  \item \textbf{Sub-CPMK-P8:} Mahasiswa mampu menerapkan \textit{ensemble} ringkas (mis.\ Random Forest) dan klastering (mis.\ K-Means) pada data yang relevan.

  \item \textbf{Sub-CPMK-P9:} Mahasiswa mampu menggunakan metrik klasifikasi/regresi, matriks konfusi, kurva ROC-AUC (jika relevan), \textit{cross-validation}, dan \textit{grid search}/\textit{random search} sederhana.

  \item \textbf{Sub-CPMK-P10:} Mahasiswa mampu menyusun ringkasan eksekutif dalam notebook (teks + visual) dan mendemonstrasikan alur \textit{end-to-end} mini (\textit{pipeline} dari data mentah ke evaluasi).
\end{itemize}

\vspace{0.5cm}

% ============================================================
% MATRIKS PEMETAAN OBE
% ============================================================
\section*{5. Matriks Pemetaan Capaian}

\textbf{Tabel 5.a. CPMK Praktikum $\rightarrow$ CPL Prodi}

\vspace{0.25cm}
{\small
\begin{tabular}{|l|p{9.5cm}|}
\hline
\textbf{CPMK} & \textbf{CPL yang dilayani} \\
\hline
CPMK-P1 & CPL01 \\
\hline
CPMK-P2 & CPL01, CPL03 \\
\hline
CPMK-P3 & CPL03 \\
\hline
CPMK-P4 & CPL01, CPL03 \\
\hline
\end{tabular}
}

\vspace{0.45cm}

\textbf{Tabel 5.b. Sub-CPMK Praktikum $\rightarrow$ CPMK Praktikum}

\vspace{0.25cm}
{\small
\begin{tabular}{|l|l|}
\hline
\textbf{Sub-CPMK} & \textbf{CPMK Praktikum} \\
\hline
Sub-CPMK-P1 & CPMK-P1 \\
\hline
Sub-CPMK-P2 & CPMK-P2 \\
\hline
Sub-CPMK-P3 & CPMK-P2 \\
\hline
Sub-CPMK-P4 & CPMK-P2 \\
\hline
Sub-CPMK-P5 & CPMK-P3 \\
\hline
Sub-CPMK-P6 & CPMK-P3 \\
\hline
Sub-CPMK-P7 & CPMK-P3 \\
\hline
Sub-CPMK-P8 & CPMK-P3 \\
\hline
Sub-CPMK-P9 & CPMK-P4 \\
\hline
Sub-CPMK-P10 & CPMK-P4 (dan CPMK-P1 untuk komunikasi masalah) \\
\hline
\end{tabular}
}

\vspace{0.45cm}

\textbf{Tabel 5.c. CPMK Praktikum $\rightarrow$ CPMK Mata Kuliah Teori \textit{Data Science} (referensi silabus teori)}

\vspace{0.25cm}
{\small
\begin{tabular}{|l|l|}
\hline
\textbf{CPMK Praktikum} & \textbf{CPMK teori terkait (\texttt{silabus\_data\_science.tex})} \\
\hline
CPMK-P1 & CPMK1 \\
\hline
CPMK-P2 & CPMK2 \\
\hline
CPMK-P3 & CPMK3 \\
\hline
CPMK-P4 & CPMK4 \\
\hline
\end{tabular}
}

\vspace{0.45cm}

\noindent\textbf{Strategi urutan pembelajaran praktikum:} Urutan minggu mengikuti alur \textbf{prosedural pipeline}: orientasi lingkungan $\rightarrow$ pemahaman data $\rightarrow$ prapemrosesan $\rightarrow$ EDA $\rightarrow$ API pemodelan $\rightarrow$ regresi $\rightarrow$ klasifikasi $\rightarrow$ \textit{ensemble}/klastering $\rightarrow$ evaluasi dan optimasi $\rightarrow$ komunikasi \textit{insight} $\rightarrow$ integrasi \textit{end-to-end}. Struktur ini membangun keterampilan prasyarat sebelum tugas proyek mini.

\vspace{0.5cm}

% ============================================================
% 6. METODE PEMBELAJARAN
% ============================================================
\section*{6. Metode Pembelajaran}

\begin{itemize}[leftmargin=*]
  \item \textbf{Briefing dan demonstrasi:} Dosen/asisten mendemonstrasikan prosedur di Jupyter sebelum praktik mandiri.
  \item \textbf{Praktikum terbimbing:} Mahasiswa menyelesaikan latihan terstruktur di notebook dengan konsultasi.
  \item \textbf{Praktik mandiri / berpasangan:} Pengerjaan modul atau tugas mingguan dengan kolaborasi terbatas sesuai aturan prodi.
  \item \textbf{Tutor sebaya / \textit{peer review}:} Pengecekan rekan terhadap kejelasan notebook dan reproduktifitas ringkas.
  \item \textbf{Proyek berbasis masalah:} Studi kasus dataset terbuka (\textit{e.g.} Kaggle, \texttt{data.go.id}) menutup alur kerja.
\end{itemize}

\vspace{0.5cm}

% ============================================================
% 7. PENGALAMAN BELAJAR MAHASISWA
% ============================================================
\section*{7. Pengalaman Belajar Mahasiswa}

\begin{itemize}[leftmargin=*]
  \item Menginstal atau mengakses lingkungan Jupyter (lokal atau Colab) dan mengunggah notebook ke sistem pengumpulan prodi.
  \item Menyimpan setiap sesi dalam berkas \texttt{.ipynb} yang dapat dijalankan ulang (\textit{restart kernel and run all}) tanpa error pada sel utama.
  \item Mengumpulkan dataset dengan atribusi sumber dan mematuhi lisensi.
  \item Mendokumentasikan langkah prapemrosesan dan keputusan analitik dalam sel Markdown.
  \item Mengekspor visualisasi yang mendukung narasi \textit{insight}.
  \item Menyusun proyek mini akhir semester berupa satu notebook atau repositori kecil sesuai panduan dosen.
\end{itemize}

\vspace{0.5cm}

% ============================================================
% 8. KRITERIA, INDIKATOR, DAN BOBOT PENILAIAN
% ============================================================
\section*{8. Kriteria, Indikator, dan Bobot Penilaian}

Mata kuliah \textbf{100\% praktikum}; tidak terdapat UTS/UAS tertulis terpisah kecuali ditetapkan lain oleh prodi. Bobot berikut total \textbf{100\%}.

\begin{longtable}{|>{\raggedright\arraybackslash}p{3cm}|>{\raggedright\arraybackslash}p{4.2cm}|>{\raggedright\arraybackslash}p{5.8cm}|c|}
\hline
\textbf{Komponen} & \textbf{Teknik Asesmen} & \textbf{Indikator / capaian utama} & \textbf{Bobot (\%)} \\
\hline
\endfirsthead
\hline
\textbf{Komponen} & \textbf{Teknik Asesmen} & \textbf{Indikator / capaian utama} & \textbf{Bobot (\%)} \\
\hline
\endhead
\hline
\endfoot

Kinerja lab / tugas notebook & Rubrik notebook per minggu atau paket modul & Sub-CPMK-P1 s.d.\ P9; CPMK-P2, P3, P4 parsial & 40 \\
\hline
Proyek mini \textit{end-to-end} & Satu notebook/repositori: data $\rightarrow$ model $\rightarrow$ evaluasi $\rightarrow$ insight & Sub-CPMK-P10; CPMK-P2, P3, P4 & 40 \\
\hline
Laporan ringkas \& presentasi & Dokumen ringkas (2--4 halaman) + presentasi singkat (5--10 menit) & CPMK-P4; komunikasi insight & 20 \\
\hline
\textbf{Total} & & & \textbf{100} \\
\hline
\end{longtable}

\textbf{Kriteria huruf (contoh):}
\begin{itemize}[leftmargin=*]
  \item A (85--100): Notebook rapi, reproduktif, analisis dan metrik tepat, insight jelas.
  \item B (70--84): Sebagian besar tujuan praktikum tercapai dengan kesalahan kecil.
  \item C (60--69): Tujuan dasar tercapai; ada kekurangan pada dokumentasi atau evaluasi.
  \item D (50--59): Hanya sebagian kecil capaian yang memadai.
  \item E ($<$50): Tidak memadai untuk kelulusan.
\end{itemize}

\textbf{Batas kelulusan:} Sesuai kebijakan fakultas/prodi (umumnya nilai akhir $\geq$ 55 atau C).

\vspace{0.5cm}

% ============================================================
% 9. EVALUASI DAN REFLEKSI PEMBELAJARAN
% ============================================================
\section*{9. Evaluasi dan Refleksi Pembelajaran}

\begin{itemize}[leftmargin=*]
  \item \textbf{Formatif:} ceklis asisten saat sesi lab, umpan balik cepat pada notebook.
  \item \textbf{Sumatif:} pengumpulan tugas notebook dan penilaian proyek mini.
  \item \textbf{Refleksi:} mahasiswa mencatat kendala teknis dan pelajaran etika data di bagian akhir notebook proyek.
  \item \textbf{Perbaikan berkelanjutan:} revisi RPS berdasarkan hasil asesmen dan survei singkat akhir semester.
\end{itemize}

\vspace{0.5cm}

% ============================================================
% 10. RENCANA PEMBELAJARAN MINGGUAN (INTI RPS PRAKTIKUM)
% ============================================================
\section*{10. Rencana Pembelajaran Mingguan}

{\footnotesize
\setlength{\tabcolsep}{3pt}%
\begin{longtable}{|>{\raggedright\arraybackslash}p{0.55cm}|>{\raggedright\arraybackslash}p{1.35cm}|>{\raggedright\arraybackslash}p{2.05cm}|>{\raggedright\arraybackslash}p{1.75cm}|>{\raggedright\arraybackslash}p{2.15cm}|>{\raggedright\arraybackslash}p{0.85cm}|>{\raggedright\arraybackslash}p{1.45cm}|>{\raggedright\arraybackslash}p{2.05cm}|}
\hline
\textbf{Mg.} & \textbf{Sub-CPMK} & \textbf{Bahan kajian / topik lab} & \textbf{Metode} & \textbf{Pengalaman belajar / artefak} & \textbf{Waktu} & \textbf{Asesmen} & \textbf{Ref.\ singkat} \\
\hline
\endfirsthead
\hline
\textbf{Mg.} & \textbf{Sub-CPMK} & \textbf{Bahan kajian / topik lab} & \textbf{Metode} & \textbf{Pengalaman belajar / artefak} & \textbf{Waktu} & \textbf{Asesmen} & \textbf{Ref.\ singkat} \\
\hline
\endhead
\hline
\endfoot

1 & P1 & Jupyter, kernel, Markdown; CRISP-DM ringkas; etika \& lisensi data. & Demo, praktik mandiri. & Notebook ``Hello Jupyter'' + refleksi etika. & 170' & Formatif & Microsoft Learn Jupyter; IBM CRISP-DM \\
\hline
2 & P2 & Pandas: baca CSV/Excel; \texttt{head}, \texttt{info}, \texttt{describe}, tipe data. & Demo, latihan. & Notebook inspeksi dataset. & 170' & Tugas modul & Kaggle Pandas; Pandas IO \\
\hline
3 & P3 & \textit{Missing}, duplikat, \textit{outlier} dasar; strategi imputasi sederhana. & PBL singkat. & Notebook cleaning. & 170' & Tugas modul & Pandas missing data; Kaggle Intermediate ML \\
\hline
4 & P3 & \textit{One-hot}, label encoding; normalisasi; fitur sederhana. & Demo, latihan. & Notebook transformasi. & 170' & Formatif & sklearn preprocessing \\
\hline
5 & P4 & EDA univariat/bivariat; distribusi, korelasi; Seaborn. & Praktik berpasangan. & Notebook EDA + minimal 4 visual. & 170' & Tugas modul & Microsoft Explore data; Seaborn tutorial \\
\hline
6 & P5 & Supervised vs unsupervised; \texttt{train\_test\_split}; alur \texttt{fit}/\texttt{predict}. & Demo, latihan. & Notebook API sklearn. & 170' & Kuis praktik singkat & sklearn tutorial; Kaggle Intro ML \\
\hline
7 & P6 & Regresi linear; interpretasi koefisien/ residual ringkas. & Praktik mandiri. & Notebook regresi + metrik RMSE/MAE. & 170' & Tugas modul & ISLR; VanderPlas handbook \\
\hline
8 & P7 & Klasifikasi: logistik, K-NN, pohon; dataset terbuka. & PBL. & Notebook klasifikasi + matriks konfusi. & 170' & Tugas modul & sklearn supervised; SciPy 2018 nb \\
\hline
9 & P8 & Random Forest singkat; perbandingan dengan baseline. & Demo, latihan. & Notebook ensemble. & 170' & Formatif & sklearn supervised \\
\hline
10 & P8 & K-Means; elbow/silhouette konsep praktik. & Praktik mandiri. & Notebook klastering. & 170' & Tugas modul & sklearn unsupervised; silhouette contoh \\
\hline
11 & P9 & Akurasi, presisi, recall, F1; ROC-AUC opsional. & Latihan terpandu. & Notebook metrik. & 170' & Rubrik sesi & sklearn model\_evaluation \\
\hline
12 & P9 & \textit{Cross-validation}; \textit{GridSearchCV} sederhana. & Demo, latihan. & Notebook tuning. & 170' & Tugas modul & sklearn CV; grid\_search \\
\hline
13 & P10 & Struktur narasi di Markdown; ringkasan untuk pemangku kepentingan. & Diskusi, praktik. & Notebook ``story'' dari minggu 7--12. & 170' & Formatif & IBM CRISP-DM \\
\hline
14 & P10 & \textit{Peer review} notebook; perbaikan reproduktifitas. & Peer review. & Revisi notebook + lembar review. & 170' & Lembar review & --- \\
\hline
15 & P10 & Proyek mini \textit{end-to-end} (individu/pasangan). & Konsultasi. & Draft proyek (\texttt{.ipynb}). & 170' & Proyek (parsial) & data.go.id; Kaggle \\
\hline
16 & P10 & Presentasi hasil proyek; umpan balik. & Presentasi. & Slides singkat + notebook final. & 170' & Presentasi \& proyek akhir & INRIA sklearn MOOC \\
\hline
\end{longtable}
}

\vspace{0.5cm}

% ============================================================
% 11. CHECKLIST SELARAS SN-DIKTI / OBE (BAGIAN F)
% ============================================================
\section*{11. Checklist Verifikasi RPS Praktikum}

\begin{enumerate}[leftmargin=*]
  \item \textbf{Standar isi:} Bahan kajian selaras IPTEKS; capaian praktik diukur lewat notebook dan rubrik, bukan sekadar daftar software.
  \item \textbf{Standar proses:} Metode menonjolkan aktivitas mahasiswa di lab; tugas dan artefak jelas.
  \item \textbf{Standar penilaian:} Bobot 100\% dengan indikator pada Sub-CPMK/CPMK; instrumen rubrik tersedia bagi dosen/asisten.
  \item \textbf{Standar sarana:} Fasilitas lab dan alternatif Colab tercantum; rencana cadangan jika server lab gangguan (opsional prodi).
  \item \textbf{Kemutakhiran:} Tanggal revisi dicatat; RPS ditinjau per semester.
  \item \textbf{Penjajaran OBE:} Setiap minggu memiliki Sub-CPMK, kegiatan, dan asesmen yang konsisten.
\end{enumerate}

\vspace{0.5cm}

% ============================================================
% 12. DAFTAR REFERENSI (sumber terbuka; selaras sumber_materi.md)
% ============================================================
\section*{12. Daftar Referensi}

\begin{enumerate}[leftmargin=*]
  \item IBM. \textit{CRISP-DM Help Overview}. \url{https://www.ibm.com/docs/en/spss-modeler/saas?topic=dm-crisp-help-overview}.

  \item IBM. \textit{SPSS Modeler CRISP-DM Guide} (PDF). \url{https://public.dhe.ibm.com/software/analytics/spss/documentation/modeler/18.0/en/ModelerCRISPDM.pdf}.

  \item Kaggle. \textit{Learn Python}. \url{https://www.kaggle.com/learn/python}.

  \item Kaggle. \textit{Learn Pandas}. \url{https://www.kaggle.com/learn/pandas}.

  \item Kaggle. \textit{Intro to Machine Learning}. \url{https://www.kaggle.com/learn/intro-to-machine-learning}.

  \item Kaggle. \textit{Intermediate Machine Learning}. \url{https://www.kaggle.com/learn/intermediate-machine-learning}.

  \item VanderPlas, J. \textit{Python Data Science Handbook}. \url{https://jakevdp.github.io/PythonDataScienceHandbook/}.

  \item James, G., dkk. \textit{An Introduction to Statistical Learning}. \url{https://www.statlearning.com/}.

  \item Microsoft Learn. \textit{Get started with Jupyter notebooks for Python}. \url{https://learn.microsoft.com/en-us/training/modules/python-create-run-jupyter-notebook/}.

  \item Microsoft Learn. \textit{Introduction to the pandas library for data science}. \url{https://learn.microsoft.com/en-us/training/modules/pandas-data-science/}.

  \item Microsoft Learn. \textit{Explore and analyze data with Python}. \url{https://learn.microsoft.com/en-us/training/modules/explore-analyze-data-with-python/}.

  \item INRIA. \textit{Machine learning in Python with scikit-learn} (materi terbuka). \url{https://inria.github.io/scikit-learn-mooc/}.

  \item VanderPlas, J. \textit{scikit-learn tutorial} (nbviewer). \url{https://nbviewer.org/github/jakevdp/sklearn_tutorial/blob/master/notebooks/Index.ipynb}.

  \item Mueller, A. SciPy 2018. \textit{Supervised Learning --- Classification} (notebook). \url{https://nbviewer.org/github/amueller/scipy-2018-sklearn/blob/master/notebooks/05.Supervised_Learning-Classification.ipynb}.

  \item scikit-learn developers. \textit{Tutorial --- Basic}. \url{https://scikit-learn.org/stable/tutorial/basic/tutorial.html}.

  \item pandas development team. \textit{User Guide: IO tools}. \url{https://pandas.pydata.org/docs/user_guide/io.html}.

  \item pandas development team. \textit{User Guide: Missing data}. \url{https://pandas.pydata.org/docs/user_guide/missing_data.html}.

  \item scikit-learn developers. \textit{Preprocessing data}. \url{https://scikit-learn.org/stable/modules/preprocessing.html}.

  \item Matplotlib development team. \textit{Matplotlib Tutorials}. \url{https://matplotlib.org/stable/tutorials/index.html}.

  \item Seaborn developers. \textit{Seaborn tutorial}. \url{https://seaborn.pydata.org/tutorial.html}.

  \item scikit-learn developers. \textit{Supervised learning}. \url{https://scikit-learn.org/stable/supervised_learning.html}.

  \item scikit-learn developers. \textit{Unsupervised learning}. \url{https://scikit-learn.org/stable/unsupervised_learning.html}.

  \item scikit-learn developers. \textit{Metrics and scoring}. \url{https://scikit-learn.org/stable/model_evaluation.html}.

  \item scikit-learn developers. \textit{Cross-validation}; \textit{Hyper-parameter tuning}. \url{https://scikit-learn.org/stable/modules/cross_validation.html} dan \url{https://scikit-learn.org/stable/modules/grid_search.html}.

  \item scikit-learn developers. Contoh \textit{silhouette analysis} K-Means. \url{https://scikit-learn.org/stable/auto_examples/cluster/plot_kmeans_silhouette_analysis.html}.

  \item Pemerintah Indonesia. \textit{Data Terbuka Indonesia}. \url{https://data.go.id/}.

  \item \textit{Perbandingan kurikulum RPS PT} (opsional dosen): contoh halaman MK \textit{Praktikum Data Science} UPN Veteran Yogyakarta. \url{https://if.upnyk.ac.id/mata-kuliah/123210411}. PDF RPS \textit{Data Wrangling} UNESA. \url{https://sindig.unesa.ac.id/rps-pdf/s1-sains-data/data-wrangling.pdf}.
\end{enumerate}

\vspace{1cm}

\begin{flushright}
\begin{tabular}{c}
Disusun oleh,\\[2cm]
\textbf{[Nama Dosen, M.Kom.]}\\
NIP. [NIP]
\end{tabular}
\end{flushright}

\end{document}
